\documentclass{beamer}

\usepackage[T2A]{fontenc}
\usepackage[utf8]{inputenc}
\usepackage[english,russian]{babel}
\input{settings.tex}   

\newcommand{\myqrframe}{
	\begin{frame}
		\centerline{Слайды доступны на Github:}
		\bigskip
		\centerline{\mygit}
		\bigskip
		\myqr
	\end{frame}
}

\subtitle{Современные методы управления роботами}
\author{Сергей Савин}
\date{2026}

\title{Теория управления: повторение}
\centering

\begin{document}
	\maketitle
	
	\begin{frame}{Содержание}
		\begin{itemize}
			\item Объекты и динамические системы
			\item Автономные и управляемые объекты
			\item Регуляторы и обратная связь по состоянию
			\item Замкнутые системы
			\item Мотивация для синтеза регуляторов и асимптотическая устойчивость
			\item Линейные стационарные системы (LTI)
			\item Устойчивость линейных систем
			\item Синтез стабилизирующего управления
			\item Методы настройки коэффициентов: размещение полюсов и LQR
			\item Линеаризация нелинейных систем
			\item Ограничения управления на основе линеаризации
		\end{itemize}
	\end{frame}
	
	\begin{frame}{Объекты и динамические системы}
		\begin{flushleft}
			
			\textbf{Объект} -- это физическая система, поведение которой изменяется во времени.
			Мы описываем её с помощью дифференциального уравнения
			\[
			\dot{x}=f(x,u),
			\]
			где \(x\) -- \textbf{состояние} объекта, а \(u\) -- его \textbf{вход}.
			
			\bigskip
			
			{\mycolor{mydarkblue}
				
				\textbf{Пример: Последовательная RC-цепь}
				
				Пусть \(x\) -- напряжение на конденсаторе, а \(u\) -- входное напряжение.
				Тогда
				\[
				\dot{x}=-\frac{1}{RC}x+\frac{1}{RC}u.
				\]
			}
		\end{flushleft}
	\end{frame}
	
	\begin{frame}{Автономные и управляемые объекты}
		\begin{flushleft}
			
			Объект, описываемый уравнением
			\[
			\dot{x}=f(x)
			\]
			не имеет внешнего управляющего входа. Такая система называется \textbf{автономной}.
			
			\bigskip
			
			Объект, описываемый уравнением \(\dot{x}=f(x,u)\), имеет вход \(u\), который может влиять
			на его поведение. Поведение автономной системы, с другой стороны, определяется только её текущим состоянием.
			
			\bigskip
			
			{\mycolor{mydarkblue}
				
				\textbf{Примеры:}
				
				\begin{itemize}
					\item  Двигатель постоянного тока -- это объект со входом: мы можем управлять им подавая напряжение.
					\item Свободный маятник -- автономеная система: его движение развивается согласно текущему состоянию, управляющего воздействия нет.
				\end{itemize}
			}
			
		\end{flushleft}
	\end{frame}
	
	\begin{frame}{Регуляторы}
		\begin{flushleft}
			
			\textbf{Регулятор} -- это правило, которое определяет, какой вход следует подавать на объект.
			
			Регулятора с обратной связью по состоянию:
			\[
			u=g(x).
			\]
			Другими словами, регулятор получает на вход текущее состояние \(x\) и определяет,
			какой управляющий сигнал \(u\) подать.
			
			\bigskip
			
			{\mycolor{mydarkblue}
				\textbf{Пример: Управление положением двигателя постоянного тока}
				
				Пропорциональный регулятор
				\[
				u=-k_p(\varphi-\varphi^\star)
				\]
				подаёт напряжение, пропорциональное ошибке положения.
			}
		\end{flushleft}
	\end{frame}
	
	\begin{frame}{Замкнутые системы}
		\begin{flushleft}
			
			Рассмотрим объект со входом:
			\[
			\dot{x}=f(x,u).
			\]
			Если мы подключим к нему регулятор \(u=g(x)\), то
			\[
			\dot{x}=f\bigl(x,g(x)\bigr).
			\]
			
			Полученная система называется \textbf{замкнутой}. Поскольку вход теперь определяется состоянием, замкнутая система является автономной.
			
			\bigskip
			
			{\mycolor{mydarkblue}
				\textbf{Пример}
				
				Водонагреватель со встроенной терморегуляцией непрерывно корректирует подачу тепла
				на основе текущей температуры -- это автономная замкнутая система.
			}
			
		\end{flushleft}
	\end{frame}
	
	\begin{frame}{Зачем мы разрабатываем регуляторы?}
		\begin{flushleft}
			
			Одна из первых целей синтеза управления -- сделать замкнутую систему
			ведущей себя желаемым образом. В частности, мы часто хотим, чтобы точка равновесия \(x^\star\) была
			\textbf{асимптотически устойчивой}.
			
			\begin{block}{Упрощённое определение}
				Если каждое решение, начинающееся достаточно близко к \(x^\star\), сходится к
				\(x^\star\), то равновесие является асимптотически устойчивым.
			\end{block}
			
			\bigskip
			
			{\mycolor{mydarkblue}
				\textbf{Пример}
				
				Самолёт, который возвращается к желаемой ориентации после малого возмущения,
				является примером устойчивого поведения в замкнутом контуре.
			}
			
		\end{flushleft}
	\end{frame}
	
	\begin{frame}{Линейные системы}
		\begin{flushleft}
			
			Особенно важным классом объектов являются \textbf{линейные стационарные (LTI)}
			системы:
			\[
			\dot{x}=Ax+Bu.
			\]
			
			Часто используемым типом регулятора является \textbf{линейный закон обратной связи по состоянию}
			\[
			u=-Kx.
			\]
			
			Подстановка регулятора в модель объекта дает замкнутую систему
			\[
			\dot{x}=(A-BK)x.
			\]
			
		\end{flushleft}
	\end{frame}
	
	\begin{frame}{Устойчивость линейных систем}
		\begin{flushleft}
			
			Рассмотрим автономную LTI-систему
			\[
			\dot{x}=Ax.
			\]
			
			Её равновесие \(x^\star=0\) является \textbf{асимптотически устойчивым} тогда и только тогда, когда
			все собственные значения \(A\) имеют строго отрицательные вещественные части:
			\[
			\Re\{\lambda_i(A)\}<0,
			\qquad \forall i.
			\]
			
			Если система асимптотически устойчива, её решения сходятся к нулю
			\textbf{с экспоненциальной скоростью}.
			
		\end{flushleft}
	\end{frame}
	
	\begin{frame}{Синтез стабилизирующего управления}
		\begin{flushleft}
			Для линейного объекта
			\[
			\dot{x}=Ax+Bu,
			\]
			и регулятора \(u=-Kx\), замкнутая система имеет вид
			\[
			\dot{x}=(A-BK)x.
			\]
			
			Таким образом, задача стабилизирующего управления заключается в следующем:
			%
	\[
	\boxed{
		\parbox{1.0\textwidth}{
			\centering
			Найти $K$ такое, что все собственные значения $A-BK$\\
			имеют отрицательные вещественные части.
		}
	}
			\]
			
			Мы можем решать эту задачу численно с помощью множества алгоритмов, например,
			\textbf{размещения полюсов} или \textbf{LQR}.
			
		\end{flushleft}
	\end{frame}
	
	\begin{frame}{Методы настройки коэффициентов}
		\begin{flushleft}
			
			Двумя распространёнными методами нахождения стабилизирующего коэффициента обратной связи \(K\) являются
			размещение полюсов и LQR.
			
			\bigskip
			
			\textbf{Размещение полюсов}:
			
			Находит \(K\) такое, что замкнутая матрица \(A-BK\) имеет в точности
			заданные собственные значения.
			
			\bigskip
			
			\textbf{Линейно-квадратичный регулятор (LQR)}:
			
			Находит \(K\), минимизирующее стоимость
			\[
			J=\int_0^\infty \left(x(t)^\top Qx(t)+u(t)^\top Ru(t)\right)dt.
			\]
			
		\end{flushleft}
	\end{frame}
	
	\begin{frame}{Линеаризация}
		\begin{flushleft}
			
			Большинство физических систем, которыми мы хотим управлять, являются \textbf{нелинейными}.
			
			\bigskip
			
			Чтобы применить линейные методы управления, мы часто аппроксимируем нелинейную
			динамику вблизи рабочей точки линейной системой. Рассмотрим:
			\[
			\dot{x}=f(x,u),
			\]
			с равновесием \((x^\star,u^\star)\).
			
			\bigskip
			
			Линеаризованная динамика для малых отклонений имеет вид
			\[
			\dot{x}=Ax+Bu,
			\qquad
			A=\left.\frac{\partial f}{\partial x}\right|_{(x^\star,u^\star)},
			\quad
			B=\left.\frac{\partial f}{\partial u}\right|_{(x^\star,u^\star)}
			\]
			
		\end{flushleft}
	\end{frame}
	
	\begin{frame}{Ограничения управления на основе линеаризации}
		\begin{flushleft}
			\begin{adjustwidth}{-0.5cm}{-0.5cm}
			
			Рассмотрим нелинейный объект
			\[
			\dot{x}=x+x^3+u.
			\]
			
			Линеаризация в окрестности \(x^\star=0\) дает
			\[
			\dot{x}=x+u.
			\]
			
			Используя размещение полюсов, мы можем выбрать
			\[
			u=-2x,
			\]
			что дает устойчивую линейную замкнутую систему с собственным значением \(-1\):
			\[
			\dot{x}=-x,
			\qquad \lambda=-1.
			\]
			
			Но замкнутая форма для исходной нелинейной системы имеет вид
			\[
			\dot{x}=-x+x^3.
			\]
			Если \(x(0)>1\), то \(\dot{x}>0\), и решение удаляется от нуля.
			
			\end{adjustwidth}
		\end{flushleft}
	\end{frame}
	
	\begin{frame}{Литература}
		
		\begin{itemize}
			
			
	\item MIT OpenCourseWare -- Feedback Control Systems: \bref{https://ocw.mit.edu/courses/16-30-feedback-control-systems-fall-2010/resources/mit16_30f10_lec05/}
{Lecture 5: State-Space Systems}

		\item MIT OpenCourseWare -- Feedback Control Systems: \bref{https://ocw.mit.edu/courses/16-30-feedback-control-systems-fall-2010/resources/mit16_30f10_lec18/}
{Lecture 18: Deterministic Linear Quadratic Regulator (LQR).}

		\item MIT OpenCourseWare -- Underactuated Robotics: \bref{https://ocw.mit.edu/courses/6-832-underactuated-robotics-spring-2009/72bc06c4dc73315bf49c28a81dc2b996_MIT6_832s09_read_ch03.pdf}
{Chapter 3: Linear Quadratic Regulator, with discussion of linearized dynamics and LQR state feedback.}
			
		\end{itemize}
		
	\end{frame}



\myqrframe





\end{document}
